\documentclass{article}
\usepackage{arxiv}
%\usepackage{amsthm}
\usepackage{amssymb}
\usepackage[utf8]{inputenc} % allow utf-8 input
\usepackage[T1]{fontenc}    % use 8-bit T1 fonts
\usepackage{mathtools}
\usepackage{hyperref}  
\usepackage{mathrsfs}

\usepackage{graphicx}
\usepackage{subcaption}
% hyperlinks
\usepackage{url}            % simple URL typesetting
\usepackage{booktabs}       % professional-quality tables
\usepackage{amsfonts}       % blackboard math symbols
\usepackage{nicefrac}       % compact symbols for 1/2, etc.
\usepackage{microtype}      % microtypography
\usepackage{lipsum}		% Can be removed after putting your text content
\usepackage{graphicx}
\usepackage{natbib}
\usepackage{doi}
\usepackage{amsmath}
\usepackage{physics}
\usepackage{wrapfig}
\usepackage{bm}
\usepackage{algorithm,algpseudocode}
\usepackage{xcolor}
\usepackage{forest}
\usepackage{tcolorbox}
\usepackage{tikz-cd}
\usepackage{enumitem}


\newenvironment{proof}{\paragraph{Proof:}}{\hfill$\square$}
\newenvironment{solution}{\paragraph{Solution:}}{\hfill$\square$}


%—– Blackboard‐bold sets —–
\newcommand{\R}{\mathbb{R}}
\newcommand{\N}{\mathbb{N}}
\newcommand{\Z}{\mathbb{Z}}
\newcommand{\Q}{\mathbb{Q}}
\newcommand{\C}{\mathbb{C}}

%—– Probability & expectation —–
\newcommand{\Prb}{\mathbb{P}}        % probability
\newcommand{\E}{\mathbb{E}}           % expectation
\newcommand{\Var}{\mathrm{Var}}       % variance
\newcommand{\Cov}{\mathrm{Cov}}       % covariance

%—– Common math operators —–
\newcommand{\diag}{\operatorname{diag}}
\newcommand{\supp}{\operatorname{supp}}
\newcommand{\argmax}{\operatorname*{arg\,max}}
\newcommand{\argmin}{\operatorname*{arg\,min}}
\newcommand{\ind}{\mathbf{1}}         % indicator

%—– Norms, inner products, absolute value —–

\newcommand{\inner}[2]{\left\langle #1, #2 \right\rangle}

%—– Calligraphic sets —–
\newcommand{\cA}{\mathcal{A}}
\newcommand{\cS}{\mathcal{S}}
\newcommand{\cT}{\mathcal{T}}
\newcommand{\cF}{\mathcal{F}}
\newcommand{\cE}{\mathcal{E}}

%—– RL / MDP notation —–
\newcommand{\SSS}{\mathcal{S}}        % state space
\newcommand{\AAA}{\mathcal{A}}        % action space
\newcommand{\PP}{\mathcal{P}}         % transition kernel
\newcommand{\RRR}{\mathcal{R}}        % reward function
\newcommand{\gammaRL}{\gamma}         % discount factor

%—– Confidence bounds —–
\newcommand{\olQ}{\overline{Q}}
\newcommand{\ulQ}{\underline{Q}}
\newcommand{\olV}{\overline{V}}
\newcommand{\ulV}{\underline{V}}

%—– Geometry / Lie groups —–
\newcommand{\SE}{\mathrm{SE}}          % special Euclidean group SE(3)
\DeclareMathOperator{\vol}{vol}        % volume

%—– Miscellaneous —–
\newcommand{\eps}{\varepsilon}
\newcommand{\dt}{\mathrm{d}t}
\newcommand{\given}{\,\bigl\vert\,}    % conditioning bar

%—– Shorthand for frequently used fractions —–
\newcommand{\half}{\tfrac12}
\newcommand{\thalf}{\tfrac32}

%—– Styling —–
\newcommand{\todo}[1]{\textcolor{red}{[TODO: #1]}}


\newcommand{\note}[1]{\begin{tcolorbox}
\textbf{Note:} #1\end{tcolorbox}}

\newtheorem{theorem}{Theorem}
\newtheorem{proposition}[theorem]{Proposition}
\newtheorem{lemma}[theorem]{Lemma}
\newtheorem{corollary}[theorem]{Corollary}
\newtheorem{definition}[theorem]{Definition}
\newtheorem{assumption}[theorem]{Assumption}
\newtheorem{remark}[theorem]{Remark}
\newtheorem{example}[theorem]{Example}
\newtheorem{conjecture}[theorem]{Conjecture}
\title{A title}
\date{} 					% Or removing it

\author{ \href{https://orcid.org/0000-0003-1700-756X}{William Chang} \\
	Department of Applied Mathematics\\
	University of California, Los Angeles\\
	Los Angeles, CA \\
	\url{https://williamc.me/} \\
	%% examples of more authors
	\And
}

% Uncomment to override  the `A preprint' in the header
\renewcommand{\headeright}{}
\renewcommand{\undertitle}{}
%\renewcommand{\shorttitle}{\textit{arXiv} Template}

%%% Add PDF metadata to help others organize their library
%%% Once the PDF is generated, you can check the metadata with
%%% $ pdfinfo template.pdf
\hypersetup{
pdftitle={},
pdfsubject={},
pdfauthor={William Chang},
pdfkeywords={},
}

\begin{document}
\maketitle
\begin{abstract}
    
\end{abstract}

% =====================================================================
%  Introduction + Preliminaries
%  Spectral-geometry view of low-frequency convergence in
%  pointmap-based 3D reconstruction.
%
%  Requires: amsmath, amssymb, amsthm, bm (optional).
%  Citation keys are placeholders -- swap for your .bib entries.
%  Suggested macros (put in preamble if not already defined):
%    \newcommand{\R}{\mathbb{R}}
%    \newcommand{\SE}{\mathrm{SE}}
%    \newcommand{\norm}[1]{\left\lVert #1 \right\rVert}
%    \DeclareMathOperator{\diag}{diag}
%    \DeclareMathOperator{\vol}{vol}
%    \DeclareMathOperator{\tr}{tr}
% =====================================================================

\section{Introduction}

Pointmap regression has reshaped multi-view 3D reconstruction. Instead of the
classical pipeline---detect correspondences, estimate poses, triangulate
depth, then fuse---models in the DUSt3R family~\cite{wang2024dust3r} learn a single
feed-forward map from an image pair $(I^1,I^2)$ to two dense \emph{pointmaps},
per-pixel 3D coordinates in a common camera frame, so that depth, relative
pose, and dense correspondence emerge together. For more than two views the
per-pair predictions must still be fused into one globally consistent scene,
and this fusion is not feed-forward: it is a per-scene optimization over
per-image depthmaps, poses, and intrinsics that minimizes pairwise pointmap
inconsistency together with flow, depth-prior, and temporal-smoothness
penalties~\cite{wang2024dust3r,zhang2025monst3r}. The network decides local geometry; the
fusion decides global consistency.

Progress along this line has been driven almost entirely by improving
\emph{correspondence}. DUSt3R is accurate on static scenes but produces noisy
pointmaps on moving objects, where its attention fails to
localize~\cite{wang2024dust3r}. zhang2025monst3r~\cite{zhang2025monst3r} augments training with dynamic
sequences but still aligns by camera pose alone. D\textsuperscript{2}USt3R~\cite{han2026enhancing},
the work we build on most directly, closes much of this gap with
\emph{static--dynamic aligned pointmaps}: it warps static pixels by the camera
pose and dynamic pixels by optical flow, yielding well-aligned ground-truth
pointmaps and markedly better reconstruction in dynamic regions. Every cue in
this progression, however, is correspondence-based---stereo matching, optical
flow, camera-induced warping---and this leaves two gaps that more
correspondence cannot fill.

First, \emph{where there is no correspondence there is no signal}. In
textureless or low-parallax regions (sky, blank walls, distant facades) and
across occlusions, neither stereo matching nor flow constrains the geometry,
so the reconstruction is hallucinated or left to drift. D\textsuperscript{2}USt3R's
machinery specifically targets \emph{dynamic} correspondence; it does nothing
for \emph{textureless static} geometry, where correspondence is equally
silent. Second, the \emph{multi-frame fusion itself} is left to a generic
alignment optimizer in which the global, large-scale degrees of
freedom---absolute scale, slowly varying depth ``bow'' across the field of
view, and cross-frame drift---are only weakly constrained and converge slowly,
the familiar behavior of the low-frequency modes in iterative geometric
solvers and bundle adjustment. D\textsuperscript{2}USt3R is itself explicit
that its two-frame model is non-generative and that extending it to video
``would require an additional global optimization technique''; that technique
is exactly what is missing.

We address both gaps with one idea: add a cue that is \emph{not}
correspondence-based. We learn a single-image map from appearance to per-pixel
surface normals (and, optionally, curvature). Appearance is informative about
local geometry even without a second view---a smooth blue region reads as a
distant, near-planar sky; shading and texture gradients reveal slope---so this
prior supplies geometric constraint precisely where correspondence is silent,
and it pins local shape directly. The difficulty is combining a single-image
cue with multi-view correspondence so the two cooperate rather than fight. We
resolve this by decomposing the reconstruction's residual in the
\emph{intrinsic spectral basis} of the geometry---the Laplace--Beltrami
operator of the reconstructed surface, realized as the graph Laplacian on the
point set~\cite{belkin2003laplacian,coifman2006diffusion,levy2006laplace}---and routing each cue to the
band where it is reliable: multi-view correspondence to the mid band; the
appearance--normal prior to the high band and to the textureless regions
correspondence cannot reach; and an explicit anchor to the single DC mode,
absolute scale, which no differential cue can supply. Because the prior's
normals are computed per depthmap in the camera frame, they are invariant to
pose and scale and therefore cannot corrupt the alignment variables: the prior
only sharpens local shape and fills holes, while correspondence and the scale
anchor own the bands it leaves alone.

This is what our method solves that the prior work does not. Relative to
D\textsuperscript{2}USt3R our contribution is complementary rather than
competing: it improves \emph{dynamic correspondence} at the per-pair, training
level, whereas we supply geometry where there is \emph{no} correspondence at
all and we give the open multi-frame fusion a concrete, geometry-aware
objective. The use of a learned monocular prior to regularize geometry where
multi-view is weak connects to monocular-prior neural reconstruction such as
yu2022monosdf~\cite{yu2022monosdf}; we differ in operating on pointmap alignment rather
than an implicit surface, and in organizing the cues by intrinsic spectral
band rather than weighting a monocular prior uniformly. The spectral viewpoint
itself draws on Laplace--Beltrami shape analysis~\cite{kac1966can,reuter2006laplace,sun2009concise}
and graph spectral filtering~\cite{hammond2011wavelets,defferrard2016convolutional}, but we use it
not as a shape \emph{descriptor} but as the scaffold that makes ``route
supervision by band'' well defined on a curved, evolving geometry---and,
because applying the operator to the coordinate maps returns the
mean-curvature normal, as the bridge to the extrinsic, bending-sensitive
curvature that the perspective projection scrambles and a flat image-plane
analysis cannot see.


% =====================================================================
%  Revised Preliminaries (trimmed) + Method
%  Organized around ONE idea: route each supervision source to the
%  intrinsic spectral band where it is reliable.
%    - mid band  : multi-view correspondence (existing alignment terms)
%    - high band : learned appearance -> normal/curvature prior  [NEW]
%    - DC        : explicit scale / gauge anchor
%  The Laplace--Beltrami / graph-Laplacian basis is what makes
%  "route by band" well defined on a curved geometry.
%
%  Requires: amsmath, amssymb, amsthm.
%  Macros assumed (see previous file): \R, \SE, \norm, \diag, \tr.
%  Add: \newtheorem{remark}{Remark}
%  Citation keys are placeholders.
% =====================================================================

% ---------------------------------------------------------------------
% Drop-in replacement for the contribution paragraph in the Intro,
% so the framing matches the Method below.
% ---------------------------------------------------------------------
\paragraph{Contributions.}
(1) We recast the global-alignment stage of pointmap reconstruction as a
\emph{band-routing} problem: its residual is a dense field whose intrinsic
spectrum, defined by the Laplace--Beltrami operator of the reconstructed
surface, separates the slow-converging global modes from the fast local ones.
(2) We observe that the two supervision sources available to the optimizer are
reliable in \emph{complementary} bands---multi-view correspondence in the mid
band, and a learned appearance-conditioned normal/curvature prior in the high
band and in textureless regions where correspondence is silent---while neither
constrains absolute scale (DC). (3) We give a single objective that routes each
source to its band, realized by \emph{local differential} operators that
require no global eigendecomposition, and we analyze its two honest failure
modes (the bas-relief low-frequency ambiguity of the appearance prior, and the
coupling between the moving surface and the basis it defines).

\section{Preliminaries}

\subsection{Pointmaps and the global-alignment objective}

Let $I^1,\dots,I^N$ be images of resolution $W\times H$. A pointmap
$X^{n,m}\in\R^{W\times H\times 3}$ gives, per pixel of image $n$, its 3D
location in the frame of camera $m$. For a pair, the network of the DUSt3R
family~\cite{wang2024dust3r,zhang2025monst3r,han2026enhancing} regresses $X^{1,1},X^{2,1}$ with
confidences $C^{v,1}$, from which depth, relative pose, and dense
correspondence follow. For $N>2$ views the per-pair predictions are fused by a
\emph{global alignment} optimization with free variables, for each image $n$, a
per-pixel log-depthmap $\theta_n$ (depth $D_n=e^{\theta_n}$), a pose
$g_n\in\SE(3)$, and a focal $f_n$ (intrinsics $K_n$); each scene-graph edge
$e=(i,j)$ carries a relative pose $g_e$, scale $\sigma_e$, and adaptor $a_e$.
The world-frame pointmap of image $n$ is
\begin{equation}
\chi_n(p)=g_n\cdot\pi^{-1}_{K_n}\!\big(p,\,D_n(p)\big),\qquad
\pi^{-1}_{K}(p,d)=d\,K^{-1}\tilde p,
\label{eq:chi}
\end{equation}
($\tilde p$ homogeneous), and the optimizer minimizes a confidence-weighted
pairwise consistency term plus flow, depth-prior, and temporal-smoothness
penalties,
\begin{equation}
\mathcal L_{\mathrm{align}}
=\!\!\sum_{e=(i,j)}\sum_{n\in\{i,j\}}\sum_p
\nu_n(p)\,\rho\!\Big(\chi_n(p)-g_e\big(\sigma_e\,a_e\!\odot\! X^{(e)}_n(p)\big)\Big),
\label{eq:align}
\end{equation}
solved by Adam for a few hundred iterations~\cite{zhang2025monst3r}. The two facts we
use are: (i) $\mathcal L_{\mathrm{align}}$ couples the variables only through
shared correspondence, so it constrains \emph{relative} depth where there is
parallax and texture but leaves absolute scale and slowly varying shape weakly
determined; and (ii) in textureless or occluded regions it provides no signal
at all.

\subsection{The intrinsic spectral basis and its curvature link}

The Fourier basis is the eigenbasis of the Laplacian; on a flat domain,
$\Delta e^{\mathrm i\langle k,x\rangle}=-\norm{k}^2 e^{\mathrm i\langle k,x\rangle}$,
so eigenvalue $=$ squared frequency. On the reconstructed surface
$\mathcal M$---with the metric induced by the 3D positions, \emph{not} the pixel
grid---the analogue is the Laplace--Beltrami operator, $\Delta_{\mathcal M}\phi_k=-\lambda_k\phi_k$,
giving a geometry-aware low/high-frequency split that is invariant to how
$\mathcal M$ projects into the image. Discretely, on the point set
$\{\chi_n(p)\}$ we use the graph Laplacian $L=\mathrm D-W$ with
$W_{ij}=\exp(-\norm{x_i-x_j}^2/\epsilon)$, whose action approximates
$\Delta_{\mathcal M}$~\cite{belkin2003laplacian}. The single property the method relies
on is that applying this operator to the coordinate map returns the
\emph{extrinsic} curvature,
\begin{equation}
\Delta_{\mathcal M}\,\mathbf x=-2H\,\mathbf n
\qquad\Longleftrightarrow\qquad
(L\chi)_n(p)\;\approx\;-\,2H_n(p)\,\mathbf n_n(p),
\label{eq:meancurv}
\end{equation}
i.e.\ twice the mean-curvature normal\footnote{Up to the sign/normalization
convention of $\Delta_{\mathcal M}$ and $H$.}. This is the bending-sensitive
quantity (intrinsic curvature is blind to bending, hence useless for the bow),
and Eq.~\eqref{eq:meancurv} is a \emph{local} operator: it needs only the
sparse $L$, never its eigenvectors.

\section{Method}

\paragraph{Overview.}
The bottleneck of Eq.~\eqref{eq:align} is spectral: local geometry (high
$\lambda$) converges in a few iterations, while the global modes (low
$\lambda$)---absolute scale, within-frame depth ``bow,'' and cross-frame
drift---decay slowly and dominate the residual, the familiar slow convergence
of the low-frequency modes in iterative geometric solvers. Our remedy splits
the reconstruction error by intrinsic spectral band, using the operator of
Eq.~\eqref{eq:meancurv}, and drives each band with the supervision that is
trustworthy there. Multi-view
correspondence is trustworthy in the mid band; we add a learned
appearance-conditioned prior for the high band and textureless fill; and
because no differential quantity carries DC information, we anchor scale
explicitly. The objective is
\begin{equation}
\boxed{\;
\mathcal L=
\underbrace{\mathcal L_{\mathrm{align}}}_{\text{mid: correspondence}}
+\,w_{\mathrm{geo}}\underbrace{\mathcal L_{\mathrm{geo}}}_{\text{high: appearance prior}}
+\,w_{\mathrm{dc}}\underbrace{\mathcal L_{\mathrm{dc}}}_{\text{DC: scale anchor}}
+\,w_{\mathrm{tmp}}\mathcal L_{\mathrm{tmp}}
+\,w_{\mathrm{dreg}}\mathcal L_{\mathrm{dreg}}\; }
\label{eq:total}
\end{equation}
with the existing temporal and depth-prior terms retained. The three boxed
terms are described below; the design principle is that they act on disjoint
spectral bands, so they cooperate rather than compete.

\subsection{The high band: a learned appearance$\to$geometry prior}

\paragraph{What we train.}
We learn a one-shot map $\mathcal N_\phi: I\mapsto(\hat N,\hat\omega)$ from a
single image to a per-pixel \emph{camera-frame} surface normal $\hat N$ and a
confidence $\hat\omega\in[0,1]$. Concretely, $\mathcal N_\phi$ is a lightweight
DPT-style head on the frozen DUSt3R encoder---whose features already encode
geometry~\cite{han2026enhancing}---supervised by normals derived from the same
ground-truth depth used to train the backbone (no new labels). Predicting the
\emph{normal} rather than raw curvature is deliberate: along the derivative
ladder $\text{depth}\!\to\!\text{normal}\!\to\!\text{curvature}$, each step is
one more high-pass ($|k|^0\!\to\!|k|^1\!\to\!|k|^2$); the normal is the rung
that is already translation- and pose-invariant (so appearance can predict it
robustly, and it cannot fight the pose variables) yet still retains slope/tilt
information, making its integration to depth a well-conditioned first-order
problem. Curvature is available as an optional second-order refinement via
Eq.~\eqref{eq:meancurv}. The ``sky is flat'' intuition is handled naturally:
in textureless regions $\mathcal N_\phi$ supplies a normal where
correspondence supplies nothing, with $\hat\omega$ down-weighting genuinely
uncertain pixels.

\paragraph{How the optimizer uses it.}
We compute the reconstruction's normals directly from each depthmap in its own
camera frame, $N_n=\mathfrak n(D_n,K_n)$, by back-projecting neighboring pixels
through Eq.~\eqref{eq:chi} and taking the normalized tangent cross product;
crucially $N_n$ depends only on $(\theta_n,f_n)$, \emph{not} on the pose $g_n$.
The high-band fidelity matches them to the prediction,
\begin{equation}
\mathcal L_{\mathrm{geo}}
=\sum_n\sum_p
\hat\omega_n(p)\,\big(1-M^n_{\mathrm{dyn}}(p)\big)
\Big[\,1-\big\langle N_n(p),\hat N_n(p)\big\rangle
+\eta\,\big\|(L\chi)_n(p)+2\hat H_n(p)\hat N_n(p)\big\|^2\Big],
\label{eq:geo}
\end{equation}
gated by the dynamic mask exactly as the flow term is. The first summand is the
first-order (normal) constraint; the optional second summand ($\eta\!\ge\!0$)
is the second-order curvature constraint of Eq.~\eqref{eq:meancurv}. Because
$\mathcal L_{\mathrm{geo}}$ is pose-free and high-pass, it pins the
\emph{within-frame} shape---a bowed depthmap has the wrong normals and is
penalized---without touching the global frame, which is precisely the
cooperation we want: it corrects the part of the bow that lives inside each
view, leaving cross-frame consistency to $\mathcal L_{\mathrm{align}}$ and
$\mathcal L_{\mathrm{tmp}}$, and scale to $\mathcal L_{\mathrm{dc}}$.

\begin{remark}[Why high-pass priors collapse the bow]
By the fundamental theorem of surface theory, fixing the (first and) second
fundamental form determines a surface up to a single global rigid motion. A
complete normal/curvature field therefore reduces the within-frame
deformation null space from an arbitrary smooth bow---infinite
dimensional---to the rigid$+$scale gauge, which we do not need to determine.
This is the precise sense in which a local, frame-free prior removes a global,
slow-converging mode.
\end{remark}

\subsection{The DC band: explicit scale anchoring}

A differential prior carries no DC: Eq.~\eqref{eq:geo} is invariant to a global
scaling $D_n\!\mapsto\! s D_n$, and so is correspondence up to the edge scales
$\sigma_e$. We therefore fix the one-dimensional scale gauge explicitly, either
by the soft anchor
$\mathcal L_{\mathrm{dc}}=\big(\log\operatorname{med}_{n,p}D_n(p)-\log s_0\big)^2$
toward a metric prior $s_0$ (e.g.\ a calibrated reference or a learned metric
depthmap on one frame), or as a hard normalization after each step. Stating
this as its own channel is not bookkeeping: it is the honest acknowledgment
that scale is the residual DC mode neither other source can supply.

\subsection{Spectral preconditioning and diagnosis (optional)}

The split of Eq.~\eqref{eq:total} rebalances which \emph{errors} are driven,
but the \emph{step size} on the depth variables is still set globally. We may
additionally amplify the low band of the depth update directly. Let $h_{\mathrm{lo}}(\lambda)$ be a smooth low-pass profile (a
spectral $\mathrm{Hann}^{1/2}$); we apply $h_{\mathrm{lo}}(L)$ to the depth
gradient via a few-term Chebyshev expansion in $L$~\cite{hammond2011wavelets,defferrard2016convolutional},
which needs only sparse matrix--vector products---no eigendecomposition---and
take a larger effective step on that component. The same operator yields a
diagnosis: projecting the residual of Eq.~\eqref{eq:align} through
$h_{\mathrm{lo}}(L)$ and its complement gives per-band convergence curves,
turning ``the global shape is still drifting'' into a measured low-mode
residual (the geometric analogue of the heat-trace/curvature readout of
spectral shape descriptors~\cite{reuter2006laplace}).

\subsection{The coupling, and why it is stable}

The graph Laplacian $L$ and the surface normals $N_n$ are computed from
$\{\chi_n\}$, which is what we are solving for, so the basis and the
constraint target both move with the iterate---the geometric face of the
problem's nonconvexity. Two things keep this benign. First, the appearance
prediction $(\hat N,\hat\omega)$ is \emph{fixed}: a one-shot function of the
input image, it is a stationary attractor toward which the moving
reconstruction is pulled, rather than a target chasing the iterate. Second,
only the cheap, local quantities ($N_n$ and the sparse $L\chi$) are recomputed
each iteration; the optional preconditioner's filter $h_{\mathrm{lo}}(L)$ is
refreshed only every few iterations, since the low modes it isolates are
exactly the ones that change slowly. We initialize $\mathcal L_{\mathrm{geo}}$
from the first iteration (the prior is valid immediately) and ramp the
optional curvature term $\eta$ and the preconditioner in only after the
warm-up used for the flow term, by which point the surface is coherent enough
for $L$ to be meaningful.

\begin{remark}[The remaining honest gap]
The appearance prior is not omniscient at low frequency: monocular shape cues
suffer the generalized bas-relief ambiguity, a low-dimensional shear/scale
freedom concentrated in the lowest modes. A learned $\mathcal N_\phi$ breaks
most of it statistically, but the residual uncertainty lives precisely in the
band we most want. Our claim is therefore not that any single channel closes
the low-frequency gap, but that the channels' reliable bands tile: the prior
pins local shape, correspondence pins relative depth across baselines, and the
anchor pins DC, so their union determines the geometry up to the global gauge
we are free to choose.
\end{remark}
\newpage

\bibliography{references}
\bibliographystyle{abbrvnat}

\end{document}